A DAG with a spine on $[0,N)$ has the edges $t\to t-1$ and further arcs $u\to p$ with $u\ge p+2$; it is an $s$-stack DAG if
its arcs split into $s$ families, each of pairwise non-crossing arcs with distinct left ends and distinct right ends. We prove
that bounded-stack DAGs reduce to constant depth: for all $s$ and $\eps>0$ there is a $D$ such that every $s$-stack DAG with a
spine has a set of at most $\eps N$ edges whose deletion leaves no path with more than $D$ arcs. Equivalently, the trajectory
cut problem $\TCP(k)$ of this series holds for every number~$k$ of walks: the walk DAG of any $k$ walks of length~$T$, with or
without pauses, has a cut of at most $\eps T$ edges and depth at most $D(k,\eps)$, independent of~$T$, where
$\log\log D(k,\eps)=O(k^2\log(k/\eps))$ for $\eps\le\frac12$. This settles $\SR(3)$, the first open case left by the eighth
paper of the series. It also shows that the deep-nesting pairs, which the series regarded as the sharpest candidates for
depth-robust walk DAGs, are cut to a depth independent of their nesting depth, so Conjectures~A$'$ and~A$''$ of the fourth and
eighth papers are false. For oblivious machines, depth reduction now holds unconditionally, and the series' reduction of a
lower bound for linear-time matrix transposition rests on the coding side alone.

The proof builds a random label from three ingredients: a decomposition of each walk into ladders of nested hops, which makes
the hops of each walk cheap on average over a random grid (Theorem~A); a display of suspended frames, which controls the
interaction between walks (view invariance); and random break bits, which keep the display bounded where exact displays are
unbounded (Lemma~D). The main theorem and every numbered result of its proof are formalized in Lean~4 with Mathlib, with no
axioms beyond \texttt{propext}, \texttt{Classical.choice} and \texttt{Quot.sound}. Not formalized are the two corollaries, for
the conjectures of the series and for oblivious machines, which combine the main theorem with results proved in the text of
earlier papers, and the remarks, which explain why simpler constructions fail or discuss the consequences.
